% Immutable cited-input supplement. Target proofs remain embargoed.
% Target: night:P4.35; tag 02X8.
% Frozen source SHA256 74497bfe327fae25c9cdcb17334b3bf42a9793a9df7c0a4404d97568f2f760f0.
% CompendiumRef is resolved against the cited S6_09 source edition, not silently identified with the current series statement.

% Exact cited input: series:II / night:P4.26
\begin{proposition}[P4.26]\hypertarget{tag-02WY}{}\label{night:P4.26}\label{stmt:P4.26}


In Part~4.2 every irreducible reduced fibre support
has nonnegative topological Euler number. For a
singular irreducible support, rational resolution
of the normalization gives strictly positive
Euler number; so does a resolution whose minimal
model is of class VII with positive second
Betti number. These assertions permit singular
normalization.

Every normal-crossing support has the
nonnegative even Euler number of P4.6.
Consequently, for a fibration all of whose
reduced supports are irreducible or have
normal crossings,
\[
 B=\mathbb P^1,\qquad c_3(Y)\in\{0,2\},
 \qquad b_3(Y)=2-c_3(Y).
\]
This last conclusion does not assume a
classification of the equality cases for
irreducible nonnormal supports.
\end{proposition}

% Exact cited input: series:II / night:P4.33
\begin{proposition}[P4.33]\hypertarget{tag-02X6}{}\label{night:P4.33}\label{stmt:P4.33}


Let $S$ be an irreducible singular support in Part~4.2, and let
$\rho:T\to N\to S$ be its minimal normalization resolution.
Then $H^2(N,\mathcal O_N)=0$. If
$p=\dim H^0(N,R^1\rho_{N*}\mathcal O_T)$ is the sum of the
geometric genera of the normalization singularities, then
\[
 p=h^1(T,\mathcal O_T)-h^1(N,\mathcal O_N).
\]
In particular a rational resolution forces every normalization
singularity to be rational and gives $\chi(\mathcal O_N)=1$.
For the scheme-theoretic image conductor $C_{\rm sch}$ we have
\[
 \chi(\mathcal O_{C_{\rm sch}})
   =-\chi(\mathcal O_N)=h^1(N,\mathcal O_N)-1.
\]

In the ruled double-horizontal case, with section $H$ of genus
$g\ge1$, a contracted $H$ gives $h^1(N,\mathcal O_N)=0$ and
$p=g$. A noncontracted $H$ gives $p=0$ and
$h^1(N,\mathcal O_N)=g$.
\end{proposition}

% Exact cited input: series:II / night:P4.8
\begin{proposition}[P4.8]\hypertarget{tag-02Y0}{}\label{night:P4.8}\label{stmt:P4.8}


For a torus fibration in P4.7 whose reduced supports have normal
crossings, an irreducible singular support with
rational normalization has Euler number two; every other irreducible
support has Euler number zero. In the zero case the supports have the
forms stated in the prompt: a smooth support has a finite unramified
torus cover; a singular irreducible support has smooth minimal
elliptic-ruled or Hopf normalization; a reducible support has smooth
minimal Hopf or non-Kähler elliptic normalization components, with
smooth elliptic double curves and finite unramified branch maps.
Thus $c_3[X]=0$ precisely when every fibre belongs to these zero cases.
\end{proposition}

% Exact cited input: series:II / night:P4R.6
\begin{proposition}[P4R.6]\hypertarget{tag-02YI}{}\label{night:P4R.6}\label{stmt:P4R.6}


Let $S$ be an irreducible singular support in
Part~4.2, let $\nu:N\to S$ be its normalization,
and let $\rho:T\to N\to S$ use the minimal
resolution of $N$. There is an effective integral
divisor $D$ on $T$ such that
\[
 \omega_T(D)\simeq\rho^*\omega_S,
 \qquad \operatorname{Supp}D
       =\rho^{-1}(\operatorname{Sing}S)_{\rm red}.
\]
In particular $K_T+D$ is torsion, by P4.12.
The strict-transform conductor has its usual
positive integral coefficients. Every connected
exceptional configuration over a singular point
of $N$ has strictly positive coefficients in $D$.
Every such configuration meets the strict-transform
conductor.
\end{proposition}

